
\documentclass[10pt,a4paper]{article}
\usepackage[left=2.5cm,right=2.5cm,top=2.5cm,bottom=2.5cm,]{geometry}
\usepackage{amsmath,amssymb,amsthm}
\newtheorem{theorem}{Theorem}[section]
\newtheorem{lemma}[theorem]{Lemma}
\newtheorem{proposition}[theorem]{Proposition}
\usepackage{setspace}
\setlength{\parindent}{0pt}
\onehalfspacing
\begin{document}
\subsubsection{dihedral group}
\begin{theorem}
    \[|D_n|=2n\]
\end{theorem}

\begin{proof}
    step 1: $|D_n|\le 2n$\\
    According to the relation of $D_n$,$D_n=\langle r,s|r^n=s^2=1,srs=r^{-1}\rangle$
    using $s^2=1\to s=s^{-1}$,using $srs=sr^{-1}\to rs=sr^{-1}$, each word can sweep to left such $r^{i_1}s^{j_1}r^{i_2}s^{j_2}\cdots$,
    until $r^is^j$ or $s^jr^i$, $i\in\{0,1,2,\cdots,n-1\}-\mathrm{Z}^+$ and $j\in\{0,1\}$ so at most have $2n$ elements. $|D_n|\le 2n$\\
    step 2: $|D_n|\ge2n$\\
    Let $R$ is rotation and $S$ is the reflection of ragular n-gon. They satisfy the same relation:
    $R^n=S^2=1,SRS=R^{-1}$
    so there is a homomorphism:
    $\phi:D_n\to\langle R,S\rangle,\phi(r)=R,\phi(s)=S$
    $\langle R,S\rangle$ is generated by $R$ and $S$. The mapping is surjection.$|D_n|\ge|\langle R,S\rangle|\ge 2n$,so $|D_n|=2n$
\end{proof}
\subsubsection{symmetric group}
For a finite group,a permutation is consists of cycles. In the other word, any permutation can write to disjoint cycle decomposition.
For a permutation, the cycle decomposition is unique, but each elements in cycle and the product of disjoin cycle commute. i.e.
\[(C_1C_2\cdots C_n)=(C_2C_3\cdots C_nC_1)=\cdots\]

\begin{lemma}
    disjoint cycle commute.
\end{lemma}

\begin{proof}
    Let $\mathrm{supp}(C)$ donated the point moves. $\forall C_i,C_j$, if $\mathrm{supp}(C_i)\cap\mathrm{supp}(C_j)=\varnothing$, 
    then $C_iC_j=C_jC_i$. $\forall x\in C_i,C_j(x)=x,C_i(x)\in\mathrm{supp}(C_i)$(cycle can permute the support set itself.) 
    so:$C_jC_i(x)=C_i(x)=C_i(x)C_j$. another case $x\in C_j$ is symmetric.
\end{proof}
 
\begin{lemma}
    \[\sigma^k=\mathrm{id.}\iff\forall i,C_i^k=\mathrm{id.}\]
\end{lemma}

\begin{proof}
    if $S_{n}^{k}=\mathrm{id.}$ and fixes $x\in\mathrm{supp}(C_i)$,for $i\ne j$, the map is the bejictive onto itself.
    the $C_i^k$ map moves points only in $\mathrm{supp}(C_i)$ and the $C_j^k$ map moves points only in $\mathrm{supp}(C_j)$. $\mathrm{supp}(C_i)\cap\mathrm{supp}(C_j)=\varnothing$
    so:$C_j^k(x)=x$\[x=\sigma^k(x)=C_1^k(x)C_2^k(x)\cdots C_i^k(x)\cdots C_j^k(x)\cdots C_n^k(x)=C_i^k(x)\]
    So $C_i^k$ fix every points in $\mathrm{supp}(C_i)$. Therefore $C_i^k=\mathrm{id.}$
\end{proof}

let $\sigma(a_i)=a_{i+1}$, if $|\sigma|=k$, $\forall m\in[1,k]$, if $\sigma^m(a_i)=a_{(i+m)\mod m}=\mathrm{id.}$, then $i+m\equiv i\mod k$, $k\mid m$, $\mathrm{ord}(\sigma)=|\sigma|$
if $\sigma^m(a_i)\ne\mathrm{id.}$, $\sigma^m$ is $\gcd(m,k)$-cycle. each size of part is $\frac{k}{\gcd(k,m)}$

\begin{lemma}
    \[\mathrm{ord}(C)=l\]
\end{lemma}

\begin{proof}
    A cycle $C=(a_1a_2\cdots a_l)$ acts $\forall j,a_j\mapsto a_{(j+1)\mod l}$,$(j+k)\equiv j\mod l\iff C^k=\mathrm{id.}\iff C^k(a_j)=a_{(j+k)\mod l}$
    so $l\mid k\iff\mathrm{ord}(C)=l$.
\end{proof}

\begin{theorem}
    $\mathrm{ord}(S_n)=\mathrm{lcm}(l_1,l_2,\cdots,l_k)$
\end{theorem}

\begin{proof}
    according to $lemma3$
\end{proof} 

more generally, let $S(\sigma)=\{k:\sigma^k=\mathrm{id.}\}$, It equals $\mathbb{Z}\triangleleft\mathrm{ord}(\sigma)\mathbb{Z}$,
acctually its the ideal of the $\mathbb{Z}$, and it can look at the cap of ideals. these two method is equival.
Then the lemma will be written $\bigcap\limits_{i=1}^{n}l_i\mathbb{Z}=\operatorname{lcm}(l_i)\mathrm{Z}$

\begin{lemma}
    $\forall a\in S$ and $p$ is prime, if $\gcd(a,p)\in\{1,p\}\land\operatorname{lcm}(a,p)=p$, then $S=\{1,p\}$
\end{lemma}

\begin{proof}
    $\forall a,a\in S$, if $a\ne p\land a\ne 1$, then $\gcd(a,p)=1,\operatorname{lcm}(a,p)=ap\ne p$,
    else if,$a=p$, then $\gcd(p,p)=p,\operatorname{lcm}(p,p)=p$,else $a=1$, then $\operatorname{lcm}(1,p)=p$,so $S=\{1,p\}$
\end{proof}

\begin{theorem}
    Given a prime $p$, an element has order-$p$ in $S_n$ iff. its cycle decomposition is a product of commuting some fixed points and $p$-cycles.
\end{theorem}

\begin{proof}
    let $C_1C_2C_3\cdots C_n$ is the joint cycle decomposition of $S_n$, 
    $\mathrm{ord}(S_n)=p\iff\operatorname{lcm}\limits_{i=1}^{|S_n|}(l_i)=p$, where $l_i=|C_i|,i\in[1,|S_n|]$
    because of $p$ is prime, $\gcd\limits_{i=1}^{|S_n|}(l_i,p)=\{1,p\}$
    According to $lemma 1$, $\{l_1,l_2,\cdots,l_i\}=\{1,p\}$.
\end{proof} 
 
Consider a group $\mathrm{Sym}(\Omega)=S_\Omega$, its joint cycle decomposition is 
$C_1C_2\cdots C_k$, $\forall i$, the generate subgroup is
$\langle C_1,C_2,\cdots,C_k\rangle\cong\prod\limits_{i=1}^{k}\langle C_i\rangle$, obviously,
$\langle C_i\rangle\cong\langle\mathbb{Z}/|C_i|\mathbb{Z},+\rangle$, $\mathbb{Z}\triangleleft|C_i|\mathbb{Z}$,
so $S_{\Omega}\cong\prod\limits_{i=1}^{k}\langle\mathbb{Z}/|C_i|\mathbb{Z},+\rangle$.

A classic assert is the $D_n$ is the automorphism of the cycle graph. Consider a cycle graph $G(V,E),V=\mathbb{Z}_n,E=E=\{\{i,i+1\}\mod n:i\in V\}$, the generate subgroup $C_n$, 
there is a bejictive $\varphi:V\to V$, $\{u,v\}\in E\iff\{\varphi(u),\varphi(v)\}\in E$.
construct:$r:i\mapsto i+1\mod|V|$ and $s:i\mapsto -i\mod|V|$, so $\langle r,s\rangle\cong\mathrm{Aut}(C_n)\subseteq\mathrm{Sym}(V)$.

Q:Normally we define the group with set and relation like $\langle S,*\rangle$ , 
almost books describe the dihedral group with generator and relation like $D_n=\langle r,s\mid r^n=s^2,srs=r^{-1}\rangle$, 
what's different?\\
A:Essentally, they both blong to $free\:group$, but the description of group consists of free group modul normal closure.
$\langle X\mid R\rangle=F(X)/\langle\langle R\rangle\rangle$. Obviously, is the quotient group. The relators are the quotirnt's id.
\begin{proposition}
    if $S_{\Omega}\cong S_{\Delta}$ and map $\theta:S_{\Omega}\to S_{\Delta}$, $\exists\varphi:S_{\Omega}\to S_{\sigma}$, $\varphi(\sigma)=\theta\circ\sigma\circ\theta^{-1}$ is surjective.
\end{proposition}
\begin{proof}
    $\theta$ is surjective, so exist $\theta^{-1}:S_{\Delta}\to S_{\Omega}$, $\forall\sigma\in S_{\Omega},\sigma:S_{\Omega}\to S_{\Omega}$.
    $\forall x\in S_{\Delta},\theta(\sigma(\theta^{-1}(x)))\in S_{\Omega}$, so $\varphi(\sigma):S_{\Delta}\to S_{\Delta}$. obviously, it's surjective.
\end{proof}
It's easy to see symmetric group only depend the cardilality of set. 
\end{document}

